% Chương 5
\chapter{KẾT LUẬN VÀ HƯỚNG PHÁT TRIỂN}
\label{Chapter5}

\section*{Tóm tắt chương}
Trong chương cuối này, chúng tôi tổng hợp những công việc đã thực hiện, đối chiếu với các mục tiêu đã đặt ra và đề xuất các hướng phát triển trong tương lai.

%----------------------------------------------------------------------------------------
%	SECTION 1
%----------------------------------------------------------------------------------------
\section{Kết luận}

Khóa luận đã nghiên cứu và xây dựng một trợ lý ảo cho hệ thống e-learning dựa trên RAG đa phương thức. Hệ thống nạp học liệu đa định dạng (tài liệu PDF, DOCX, PPTX và video bài giảng) vào một kho tri thức gồm hai chỉ mục vector và một chỉ mục từ vựng, truy xuất bằng bộ truy xuất lai hợp nhất bằng RRF có lọc theo môn học, và sinh câu trả lời kèm trích dẫn có vị trí cụ thể, với cơ chế từ chối theo ngưỡng và API truyền phát câu trả lời (\textbf{Chương \ref{Chapter3}}). Chúng tôi cũng xây dựng một quy trình đánh giá truy xuất không rò rỉ dữ liệu, trong đó truy vấn và nhãn chuẩn được tách riêng và các câu hỏi không thể trả lời hoặc có bằng chứng chưa được lập chỉ mục được xử lý riêng.

\bigbreak\noindent
Đối chiếu với bốn mục tiêu của khóa luận (\textbf{Mục \ref{sec:muc-tieu}}): mục tiêu tổ chức kho tri thức thống nhất và mục tiêu xây dựng bộ truy xuất kết hợp nhiều nguồn đã được hiện thực; mục tiêu sinh câu trả lời có căn cứ đã được hiện thực nhưng chưa được đánh giá bằng thực nghiệm; mục tiêu đánh giá truy xuất đã được thực hiện trên bộ 400 câu hỏi, với 325 câu hợp lệ và độ phủ của kho học liệu là 85,53\%.

\bigbreak\noindent
Kết quả thực nghiệm cần được đọc thận trọng (\textbf{Chương \ref{Chapter4}}). Ở cả hai hệ thống, chỉ khoảng 40\% số câu hợp lệ có ít nhất một kết quả đúng trong 10 kết quả đầu. Hệ thống đề xuất, thêm truy xuất thưa và truy xuất trên ảnh vào hệ thống cơ sở chỉ dùng văn bản, không tạo ra khác biệt có ý nghĩa thống kê ở $K=1$, $K=5$ hay $K=10$ khi lấy mẫu theo truy vấn; ở $K=5$ điểm của hệ thống đề xuất chỉ có xu hướng thấp hơn. Vì vậy, với cấu hình hiện tại (hợp nhất bằng RRF không trọng số) và kho học liệu chưa đầy đủ, chúng tôi chưa có bằng chứng rằng việc bổ sung các nguồn này cải thiện chất lượng truy xuất. Nguyên nhân chưa được kiểm chứng, và các thí nghiệm loại trừ cần thiết để làm rõ vẫn là công việc tiếp theo. Ngoài ra, khóa luận chưa đánh giá chất lượng của câu trả lời được sinh.

\bigbreak\noindent
Đóng góp của khóa luận vì thế nằm ở ba điểm: một hệ thống chạy cục bộ, được thiết kế hướng tới thiết bị có bộ nhớ hạn chế, cho phép truy xuất và trích dẫn trên học liệu đa định dạng; một quy trình đánh giá truy xuất có thể tái lập và không rò rỉ dữ liệu; và một bộ kết quả ban đầu được báo cáo trung thực, kèm các hạn chế và trạng thái của những thí nghiệm còn lại.

%-----------------------------------
%	SECTION 2
%-----------------------------------
\section{Hướng phát triển}

Dựa trên kết quả đạt được và các hạn chế đã nêu (\textbf{Mục \ref{sec:han-che}}), chúng tôi đề xuất các hướng phát triển sau.

\bigbreak\noindent
\textbf{Hoàn thiện đánh giá.} Trước hết cần lập chỉ mục đầy đủ kho học liệu rồi chạy lại hai hệ thống, sau đó thực hiện các thí nghiệm loại trừ để xác định thành phần nào ảnh hưởng đến kết quả, và mở rộng sang đánh giá đầu-cuối. Đối với phần sinh, khung RAGAs cung cấp các chỉ số cho khả năng khai thác ngữ cảnh một cách trung thực và chất lượng của phần sinh mà không cần nhãn chuẩn do con người gán \cite{es2024ragas}, và mã nguồn của hệ thống đã có sẵn chức năng chấm điểm theo hướng này. Các so sánh giữa hệ thống nên được kiểm định thống kê, ví dụ bằng bootstrap (lấy mẫu theo truy vấn khi câu hỏi bị lặp lại) hoặc các kiểm định đã được so sánh trong đánh giá truy xuất thông tin \cite{smucker2007comparison}; đồng thời cần người thứ hai rà soát nhãn chuẩn và cần xử lý các câu hỏi trùng nội dung trong bộ câu hỏi.

\bigbreak\noindent
\textbf{Hợp nhất thích nghi.} RRF chỉ dùng thứ hạng và coi mọi danh sách như nhau \cite{cormack2009rrf}, trong khi khảo sát về RAG đa phương thức xem hợp nhất là một thành phần riêng của hệ thống \cite{abootorabi2025ask}. Một hướng phát triển là gán trọng số cho từng nguồn theo loại câu hỏi và theo độ tin cậy của nguồn (ví dụ độ tin cậy của phiên âm và của nhận dạng ký tự), điều này đòi hỏi bổ sung siêu dữ liệu độ tin cậy cho từng đơn vị tri thức. Hướng này chưa được hiện thực.

\bigbreak\noindent
\textbf{Xếp hạng lại và biểu diễn đa phương thức.} Việc xếp hạng lại bằng mô hình xét đồng thời câu hỏi và đoạn văn đã được áp dụng cho truy xuất đoạn văn \cite{nogueira2019passage}; hệ thống đã có chỗ tích hợp cho Qwen3-Reranker nhưng mặc định tắt và chưa được đánh giá \cite{zhang2025qwen3embedding}. Họ Qwen3-VL cũng có mô hình xếp hạng lại đa phương thức đi kèm mô hình biểu diễn \cite{li2026qwen3vlembedding}, có thể được thử nghiệm cho các kết quả dạng ảnh.

\bigbreak\noindent
\textbf{Truy xuất trực tiếp trên ảnh trang và trên video.} Các công trình truy xuất trực tiếp trên ảnh trang tài liệu \cite{faysse2025colpali, yu2025visrag, tanaka2025vdocrag} và truy xuất ở mức video \cite{jeong2025videorag} gợi ý các hướng khai thác thị giác đầy đủ hơn so với việc chỉ đưa khung hình đã chọn vào chỉ mục và chỉ đưa ảnh trang vào bước sinh.

\bigbreak\noindent
\textbf{Siêu dữ liệu học tập và đồ thị tri thức.} Lược đồ siêu dữ liệu hiện chỉ có môn học, loại nội dung và vị trí. Có thể mở rộng lược đồ, tham khảo tiêu chuẩn LOM \cite{ieee2020lom}, bằng các trường như chương, tuần, mục tiêu học tập, độ khó và phiên bản học liệu (danh sách trường này là đề xuất của chúng tôi, không phải danh sách nêu trong tiêu chuẩn), và liên kết các nguồn về cùng một khái niệm bằng một đồ thị tri thức; một cách làm đã được đề xuất là dựng đồ thị thực thể từ tài liệu rồi tóm tắt sẵn theo cộng đồng các thực thể gần nhau để trả lời các câu hỏi tổng quát \cite{edge2024graphrag}.

\bigbreak\noindent
\textbf{Hỗ trợ bài tập và kiểm soát đáp án.} Bài tập hiện được xử lý như văn bản thông thường. Một trợ lý học tập cần cân nhắc việc không tiết lộ lời giải: CodeAid là một ví dụ về trợ lý lập trình dựa trên LLM được thiết kế để đưa ra phản hồi đúng kỹ thuật mà không tiết lộ mã lời giải \cite{kazemitabaar2024codeaid}. Việc tách đề, gợi ý, lời giải và mức độ khó, cùng cơ chế phân quyền theo vai trò người dùng và trạng thái bài tập, là hướng mở rộng cho hệ thống của chúng tôi.

\bigbreak\noindent
\textbf{Triển khai.} Dịch vụ hiện chạy trong một tiến trình và xử lý tuần tự từng yêu cầu suy luận; để phục vụ nhiều người học đồng thời cần tách tiến trình suy luận khỏi tiến trình phục vụ và quản lý hàng đợi yêu cầu.
